\documentclass[10pt,a4paper]{article}

\usepackage[utf8]{inputenc}
\usepackage[T1]{fontenc}
\usepackage[margin=1.6cm]{geometry}
\usepackage{enumitem}
\usepackage{titlesec}
\usepackage{xcolor}
\usepackage[hidelinks]{hyperref}

\definecolor{accent}{HTML}{1F4E79}

\pagestyle{empty}
\setlength{\parindent}{0pt}
\setlength{\parskip}{0pt}

% Títulos de seção
\titleformat{\section}{\large\bfseries\color{accent}\MakeUppercase}{}{0em}{}[\vspace{-0.6em}\rule{\linewidth}{0.5pt}]
\titlespacing*{\section}{0pt}{10pt}{6pt}

% Lista compacta
\setlist[itemize]{leftmargin=1.2em, itemsep=2pt, topsep=3pt, parsep=0pt}

% Entrada: título/local à esquerda, data à direita
\newcommand{\entry}[4]{%
  \textbf{#1} \hfill #3 \\
  \textit{#2} \hfill \textit{#4}\vspace{2pt}%
}

\begin{document}

% ---------- Cabeçalho ----------
\begin{center}
  {\Huge\bfseries Lorena Moreira de Nadai}\\[4pt]
  {\large\color{accent} Desenvolvimento de Software \& Engenharia de Dados}\\[6pt]
  \href{mailto:lorenadenadai679@gmail.com}{lorenadenadai679@gmail.com} \quad|\quad
  (27) 99739-9041 \quad|\quad Cariacica -- ES\\[2pt]
  \href{http://www.linkedin.com/in/lorena-denadai-17a39037b}{LinkedIn} \quad|\quad
  \href{https://github.com/denadaiLorena}{GitHub}
\end{center}

% ---------- Resumo ----------
\section{Resumo Profissional}
Estudante de Ciência da Computação (6º período) com sólida experiência em modelagem e
processamento de dados, criação de rotinas automatizadas, SQL e desenvolvimento de software.
Atualmente atua no setor de tecnologia do IEMA, sendo responsável pela modelagem de bancos
de dados, construção de pipelines batch e dashboards analíticos para acompanhamento de
indicadores institucionais. Perfil analítico, focado em organização, qualidade de dados e
soluções escaláveis.

% ---------- Experiência ----------
\section{Experiência Profissional}
\entry{Estagiária de Desenvolvimento de Software}
      {Instituto Estadual de Meio Ambiente e Recursos Hídricos (IEMA)}
      {Nov 2025 -- atual}{}
\begin{itemize}
  \item Desenvolvimento autônomo (full-stack) de sistema interno para acompanhamento de
        processos complexos, com implementação de motor de cálculo automatizado para
        estimativa de custos, validação eletrônica da documentação de solicitações e módulo
        de contabilidade e controle rigoroso do estoque consumido pelo setor interno,
        utilizando PHP, JavaScript e PostgreSQL.
\end{itemize}

% ---------- Formação ----------
\section{Formação Acadêmica}
\entry{Bacharelado em Ciência da Computação}
      {FAESA -- Centro Universitário}
      {Conclusão prevista: Dez 2027}{}
\vspace{4pt}

\entry{Montagem e Manutenção de Computadores}
      {Senai Vitória}
      {Conclusão: 2021}{}

% ---------- Projetos ----------
\section{Projetos Destacados}
\textbf{Dashboard Analítico de Indicadores (GitLab)}
\begin{itemize}
  \item Desenvolvimento de aplicação full-stack de dashboard analítico para indicadores do
        GitLab, utilizando Quasar e Flask, integrada a uma rotina batch em C\# para
        processamento automatizado de dados e PostgreSQL.
\end{itemize}

\vspace{4pt}
\textbf{Sistema de Análise e Otimização de Dados da COVID-19}
\begin{itemize}
  \item Tratamento e análise de grandes volumes de microdados de saúde utilizando Python e
        Pandas.
  \item Implementação de técnicas de \textit{chunking} e do formato Parquet, otimizando o
        carregamento e a performance no processamento de milhões de registros.
\end{itemize}

\vspace{4pt}
\textbf{Mapa de Acessibilidade (Projeto Acadêmico)}
\begin{itemize}
  \item Desenvolvimento de aplicativo baseado em geolocalização utilizando Java e a API
        osmdroid para identificar pontos urbanos acessíveis.
\end{itemize}

% ---------- Competências ----------
\section{Competências Técnicas}
\begin{itemize}
  \item[] \textbf{Linguagens \& Frameworks:} Java (Spring Boot), PHP, Python (Pandas),
          JavaScript, HTML, CSS, Quasar, Angular.
  \item[] \textbf{Bancos de Dados:} PostgreSQL, MongoDB.
  \item[] \textbf{Ferramentas \& Conceitos:} Git, GitHub, APIs RESTful, Modelagem de Dados,
          Otimização, Data Warehouse (DW).
\end{itemize}

\end{document}
